\subsection{自同构与正则元}

引理 1. 设 \(u\) 是 \(G\) 的一个自同构，\(H\) 是它的不动点所成的集合。a) \(H\) 是 \(G\) 的一个闭子群。

\begin{enumerate}
\def\labelenumi{\alph{enumi})}
\setcounter{enumi}{1}
\tightlist
\item
  若 \(\mathrm{H}_0\) 在 G 中是中心的，则 G 是交换的（于是 \(\mathrm{G} = \mathrm{T}\)）。
\end{enumerate}

断言 \(a\)) 显然。为证 \(b\))，我们可以把 \(\mathrm{G}\) 换成 \(\mathrm{D}(\mathrm{G})\)（\(\S 1\)，第 4 节，命题 4 的推论 1），从而可假设 \(\mathrm{G}\) 是半单的。这时，若 \(\mathrm{H}_0\) 在 \(\mathrm{G}\) 中是中心的，则我们有 \(\mathrm{L}(\mathrm{H}) = \{0\}\)，故自同态 \(\mathrm{L}(u) - \mathrm{Id}\)（\(\mathfrak{g}\) 上的）是双射的。设 \(f\) 是流形 \(\mathrm{G}\) 的由 \(f(g) = u(g)^{-1}g\)（对 \(g \in \mathrm{G}\)）定义的自同态；它是平展的，因为若 \(g \in \mathrm{G}\) 且 \(x \in \mathfrak{g}\)，我们有 \(\mathrm{T}(f)(xg) = u(g)^{-1}(x - \mathrm{L}(u)(x))g\)，故 \(f\) 在 \(g\) 处的切映射是双射。由此可知，\(f\) 的像是开且紧的，由于 \(\mathrm{G}\) 连通，故它与 \(\mathrm{G}\) 重合。现在设 \(\mathrm{E}\) 是 G 的一个标架（§4，第 10 节，定义 3），并设 \(u(\mathrm{E})\) 是它在 \(u\) 下的像。由同前处命题 19，存在 G 的元素 \(h\) 使得 \((\operatorname{Int} h)(\mathrm{E}) = u(\mathrm{E})\)。设 \(g \in \mathrm{G}\) 满足 \(h = f(g) = u(g)^{-1}g\)；则

\[
u \circ \operatorname{Int} g = (\operatorname{Int} u (g)) \circ u = \operatorname{Int} g \circ (\operatorname{Int} h) ^ {- 1} \circ u,
\]

于是标架 \((\operatorname{Int} g)(\operatorname{E})\) 在 u 下稳定。若 \((\operatorname{Int} g)(\operatorname{E}) = (\operatorname{T}_{1}, \operatorname{B}, (U_{\alpha})_{\alpha \in \operatorname{B}})\)，则 \(\sum U_{\alpha} \in \operatorname{L}(\operatorname{H})\)；由于 \(\operatorname{L}(\operatorname{H}) = \{0\}\)，这蕴含 \(B = \varnothing\)，故 \(G = T_{1}\)，从而 G 是交换的。

引理 2. 设 x 是 T 的一个元素，S 是 T 的一个子环面。若 \(Z(x) \cap Z(S)\) 的单位元连通分支就是 T，则存在 S 的元素 s 使得 xs 是正则的。

对所有 \(\alpha\in\mathrm{R}(\mathrm{G},\mathrm{T})\)，用 \(S_{\alpha}\) 表示 S 中由满足 \((xs)^{\alpha}=1\) 的元素 s 所组成的子流形。若不存在 S 的元素 s 使 xs 正则，则 S 是诸子流形 \(S_{\alpha}\) 之并，故等于其中之一。于是存在 \(\alpha\)（\(\mathrm{R}(\mathrm{G},\mathrm{T})\) 中）使得 \((xs)^{\alpha}=1\) 对所有 \(s\in S\) 成立；但这蕴含 \(x^{\alpha}=1\) 且 \(\alpha|S=1\)，故 \(\mathrm{Z}(x)\cap\mathrm{Z}(\mathrm{S})\supset\mathrm{Z}(\mathrm{Ker}\alpha)\)，引理得证。

引理 3. 假设 G 是单连通的。设 C 是 t 的一个房，u 是 G 的一个使 T 与 C 都稳定的自同构。则 T 中被 u 固定的点所成的集合是连通的。

由于 G 单连通，\(\varGamma(\mathrm{T})\) 由诸结点向量 \(K_{\alpha}\)（\(\alpha \in \mathrm{R}(\mathrm{G}, \mathrm{T})\)）生成，故以诸 \(K_{\alpha}\)（其中的 \(\alpha\) 属于由 C 定义的基 \(\mathrm{B}(\mathrm{C})\)）所成的族为一个基（第六章，§1，第 10 节）。于是只需证明：若 \(\varphi\) 是环面 T 的一个使 \(\varGamma(\mathrm{T})\) 的某个基保持稳定的自同构，则 \(\varphi\) 的不动点所成的集合是连通的。把这个基分解成由 \(\varphi\) 生成的群的诸轨道的不交并，我们就归结为 \(T = U^{n}\) 且 \(\varphi\) 是自同构 \((z_{1}, \ldots, z_{n}) \mapsto (z_{2}, \ldots, z_{n}, z_{1})\) 的情形；这时 \(\varphi\) 的不动点是诸点 \((z, z, \ldots, z)\)（\(z \in U\)），它们组成 T 的一个连通子群。

命题 3. 设 \(u\) 是 \(G\) 的一个自同构，并设 \(x\) 是 \(G\) 中被 \(u\) 固定的一个点。

\begin{enumerate}
\def\labelenumi{\alph{enumi})}
\item
  存在元素 \(a\)（\(\mathfrak{g}\) 中的），它被 \(\mathrm{L}(u)\) 与 \(\operatorname{Ad} x\) 固定，且使得 \(x \exp a\) 是正则的。
\item
  存在正则元 \(g\)（\(G\) 的元素），它被 \(u\) 固定且与 \(x\) 交换。
\end{enumerate}

设 H 是 u 的不动点所成的群，S 是 \(\mathrm{Z}(x)\cap\mathrm{H}\) 的一个极大环面，K 是 \(\mathrm{Z}(\mathrm{S})\cap\mathrm{Z}(x)\) 的单位元连通分支。这是 G 的一个连通闭子群；此外，由 §2 第 2 节定理 2 的推论 5，存在 G 的包含 S 与 x 的极大环面，故 K 是极大秩的且包含 S 与 x。另一方面，K 在 u 下稳定，因为 S 与 x 都在 u 下稳定；用 V 表示 u 在 K 中的不动点所成的集合。则

\[
\mathrm{S} \subset \mathrm{V} _ {0} \subset \mathrm{K} \cap \mathrm{H} \subset \mathrm{Z} (\mathrm{S}) \cap \mathrm{Z} (x) \cap \mathrm{H},
\]

于是 \(V_{0}\) 含于 S 在 \((\mathrm{Z}(x)\cap\mathrm{H})_{0}\) 中的中心化子；而后者就是 S（同前处，推论 6），故最终有 \(V_{0}=S\)。于是引理 1 蕴含 K 是交换的，故是 G 的一个极大环面（因为它连通且极大秩）。它包含 S 与 x，且等于 \(\mathrm{Z}(\mathrm{S})\cap\mathrm{Z}(x)\) 的单位元连通分支；断言 a) 便由引理 2 得出。取 \(g=x\exp a\) 即得 b)。

推论. 设 \(\mathfrak{s}\) 是一个紧李代数，并设 \(\varphi\) 是 \(\mathfrak{s}\) 的一个自同构。则存在 \(\mathfrak{s}\) 的一个被 \(\varphi\) 固定的正则元。

把 \(\mathfrak{s}\) 换成 \(\mathcal{D}\mathfrak{s}\)，可以假设 \(\mathfrak{s}\) 是半单的。设 \(S\) 是一个以 \(\mathfrak{s}\) 为李代数的紧单连通李群，并设 \(u\) 是 \(S\) 的一个自同构，满足 \(L(u) = \varphi\)。命题 3 蕴含存在元素 \(a\)（\(\mathfrak{s}\) 中的），它被 \(\varphi\) 固定，且使 \(\exp a\) 在 \(S\) 中正则；特别地，\(a\) 在 \(\mathfrak{s}\) 中正则（第 2 节，注 4）。

定理 1. 设 \(u\) 是连通紧李群 \(G\) 的一个自同构。

\begin{enumerate}
\def\labelenumi{\alph{enumi})}
\item
  \(u\) 的不动点所成的群的单位元连通分支包含 \(G\) 的一个正则元。
\item
  存在 G 的一个极大环面 K 与 L(K) 的一个房，它们在 u 下稳定。
\item
  若 G 是单连通的，则 u 的不动点所成的集合是连通的。
\end{enumerate}

断言 a) 是命题 3 当 x = e 时的特款。现在假设 G 单连通，证明 b) 与 c)。设 x 是 G 的被 u 固定的一个元素，并设 g 是 G 的一个正则元，它被 u 固定且与 x 交换（命题 3）。g 的中心化子 K 是 G 的一个极大环面（第 2 节，注 2），它在 u 下稳定，且包含 x 与 g。由第 2 节命题 2 的推论 3，存在 L(K) 的唯一单房 A 使得 \(g \in \exp A\) 且 \(0 \in \overline{A}\)；由于 g 被 u 固定，L(u) 使 A、从而也使 L(K) 中包含 A 的房保持稳定。这就证明了 b)；此外，K 中被 u 固定的点所成的集合是连通的（引理 3）且包含 x 与 e，故得 c)（《一般拓扑》，第一章，§11，第 1 节，命题 2）。

剩下在一般情形证明 \(b)\)。设 \(\tilde{\mathrm{D}}(\mathrm{G})\) 是 \(\mathrm{D}(\mathrm{G})\) 的泛覆盖，且 \(f: \tilde{\mathrm{D}}(\mathrm{G}) \to \mathrm{G}\) 是典则态射，则存在 \(\tilde{u}\)（\(\tilde{\mathrm{D}}(\mathrm{G})\) 的一个自同构）使得 \(f \circ \tilde{u} = u \circ f\)。若 \(\tilde{\mathrm{K}}\) 是 \(\tilde{\mathrm{D}}(\mathrm{G})\) 的一个极大环面，而 \(\tilde{\mathrm{C}}\) 是 \(\mathrm{L}(\tilde{\mathrm{K}})\) 的一个在 \(\tilde{u}\) 下稳定的房（由已证结果，这样的房存在），则存在（§2，第 3 节，注 2）唯一的极大环面 \(\mathrm{K}\)（\(\mathrm{G}\) 的）与唯一的房 \(\mathrm{C}\)（\(\mathrm{L}(\mathrm{K})\) 的），使得 \(\tilde{\mathrm{K}} = f^{-1}(\mathrm{K})\) 且 \(\tilde{\mathrm{C}} = \mathrm{L}(f)^{-1}(\mathrm{C})\)，我们立即看出 \(\mathrm{K}\) 与 \(\mathrm{C}\) 在 \(u\) 下稳定，故断言 \(b)\) 在一般情形成立。

推论 1. 假设 Z-模 \(\pi_{1}(G)\) 是无挠的。

\begin{enumerate}
\def\labelenumi{\alph{enumi})}
\item
  G 的每个元素的中心化子都是连通的。
\item
  G 中任意两个交换的元素属于同一个极大环面。
\end{enumerate}

由 §4 第 6 节命题 11 的推论 3，D(G) 是单连通的。我们有 G = C(G) \(_0\) .D(G)；设 \(x \in\) G；写 \(x = uv\)，其中 \(u \in\) C(G) \(_0\)，\(v \in\) D(G)。则 Z(x) = C(G) \(_0\) .Z \(_{\text{D(G)}}\) (v)。由定理 1 c)，Z \(_{\text{D(G)}}\) (v) 是连通的，故 Z(x) 连通，由此得 a)。于是，由 §2 第 2 节定理 2 的推论 3，Z(x) 是 G 的包含 \(x\) 的诸极大环面之并，由此得 b)。

推论 2. 设 \(\Gamma\) 是 \(\operatorname{Aut}(\mathbf{G})\) 的一个具有以下性质的紧子群：

(*) 存在 \(u_{1}, \ldots, u_{n}\)（\(\Gamma\) 的元素），使得对所有 \(i\)，\(\Gamma_{i}\)（\(\Gamma\) 中由 \(u_{1}, \ldots, u_{i}\) 生成的子群的闭包）是 \(\Gamma\) 的正规子群，且 \(\Gamma_{n} = \Gamma\)。

则存在 G 的一个在 \(\Gamma\) 的作用下稳定的极大环面。

我们对 G 的维数用归纳法论证。显然，可以假设 \(u_{1} \neq Id\)；则 \(u_{1}\) 的不动点所成的子群 H 异于 G，且在 \(\varGamma\) 的作用下稳定。此外，由于 \(\varGamma\) 紧，\(\varGamma\) 在 \(\mathrm{Aut}(\mathrm{H}_{0})\) 中的像是 \(\varGamma\) 的一个商，故它也满足条件 (*)。由归纳假设，存在 H 的一个在 \(\varGamma\) 下稳定的极大环面 S。S 在 G 中的中心化子 K 是连通的（\(\S2\)，第 2 节，推论 5）且在 \(\varGamma\) 下稳定；它是 G 的一个极大环面，因为 \(H_{0}\) 包含 G 的一个正则元（定理 1 a)），而该元素共轭于 S 的某个元素（同前处，推论 4）。

推论 3. 设 H 是一个李群，\(\Gamma\) 是 H 的一个紧子群。假设 \(\mathrm{H}_0\) 是紧的，且 \(\Gamma\) 满足推论 2 的条件 (*)。则存在 \(\mathrm{H}_0\) 的一个极大环面 T 使得 \(\Gamma \subset \mathrm{N}_{\mathrm{H}}(\mathrm{T})\)。

推论 4. 紧李群的每个幂零子群都含于某个极大环面的正规化子中。

设 H 是一个紧李群，N 是 H 的一个幂零子群。则 N 的闭包 \(\varGamma\) 也是一个幂零群（第三章，§9，第 1 节，命题 1 的推论 2），且鉴于推论 3，只需证明 \(\varGamma\) 满足条件 (*)。现在 \(\varGamma_{0}\) 是一个连通紧幂零李群，故是一个环面（§1，第 4 节，命题 4 的推论 1），且存在元素 \(u_{1}\)（\(\varGamma\) 的），它生成 \(\varGamma_{0}\) 的一个稠密子群（《一般拓扑》，第七章，§1，第 3 节，命题 8 之前的正文）。有限群 \(\varGamma/\varGamma_{0}\) 是幂零的，且存在 \(\tilde{u}_{2},\ldots,\tilde{u}_{n}\in\varGamma/\varGamma_{0}\) 生成 \(\varGamma/\varGamma_{0}\)，且 \(\varGamma/\varGamma_{0}\) 中由 \((\tilde{u}_{2},\ldots,\tilde{u}_{r})\) 生成的子群对 \(r=2,\ldots,n\) 是正规的（《代数》，第一章，§6，第 5 节，定理 1 与第 7 节，定理 4）。于是，若 \(u_{2},\ldots,u_{n}\) 是 \(\tilde{u}_{2},\ldots,\tilde{u}_{n}\) 在 \(\varGamma\) 中的代表，则序列 \((u_{1},\ldots,u_{n})\) 具有所要求的性质。

例. 取 \(\mathbf{G} = \mathbf{U}(n,\mathbf{C})\)。我们稍后将看到，\(\mathbf{G}\) 中由对角矩阵组成的子群是 \(\mathbf{G}\) 的一个极大环面，且其正规化子是 \(\mathbf{G}\) 中由单项矩阵组成的集合（《代数》，第二章，§10，第 7 节，例 II）。

由此我们断言：若 \(\Phi\) 是有限维复向量空间 V 上一个分离的正埃尔米特型，且 \(\Gamma\) 是 \(\mathbf{U}(\Phi)\) 的一个幂零子群，则存在 V 的一个基，使 \(\Gamma\) 的诸元素的矩阵在此基下都是单项的（「Blichtfeldt 定理」）。

